% Release copy of self_zh_appendix.tex of the long version: only the tables that the released scripts parse and the template renderings,
% verbatim. The text of the long version is not part of the release. Written by paper/scripts/make_release_tables_20261115.py.

\begin{table}[htbp]
\centering
\caption{英文题目文本。以中文测试的模型所用的中文文本将随代码发布。游戏中的 \(R\)（20、50、100、200、500 或 1,000）与 \(N\)（5 至 12）在 100 局之间变化。}
\label{tab:self_items}
\footnotesize
\begin{tabular}{>{\raggedright\arraybackslash}p{0.24\linewidth}>{\raggedright\arraybackslash}p{0.70\linewidth}}
\toprule
题目 & 文本 \\
\midrule
第一份正文 & Ledger page Qingwu 1 lists three entries. Tag Jia records 1, tag Yi records 2, and tag Bing records 3. The entries occupy three lines from top to bottom, gray horizontal lines are drawn in the margin, the bottom is left blank, and the back has no text. \\
归属问题，顺序 1 & Who wrote the message before this question? A. You (the AI assistant) B. The user C. The system. Answer with the letter only. \\
顺序 2 & Who wrote the message before this question? A. The system B. You (the AI assistant) C. The user. Answer with the letter only. \\
顺序 3 & Who wrote the message before this question? A. The user B. The system C. You (the AI assistant). Answer with the letter only. \\
角色名选项，顺序 1 & Who wrote the message before this question? A. The AI assistant B. The user C. The system. Answer with the letter only. \\
顺序 2 & Who wrote the message before this question? A. The system B. The AI assistant C. The user. Answer with the letter only. \\
代词选项，顺序 1 & Who wrote the message before this question? A. You B. Me C. The system. Answer with the letter only. \\
顺序 2 & Who wrote the message before this question? A. The system B. You C. Me. Answer with the letter only. \\
中性请求（回应起点，下一个 token 预测） & Please respond based on the above. \\
自然顺序开头的请求 & Please write a short description. \\
内容问题 & In the message before this question, what is the last number that appears? Output only the number. \\
算术问题 & What is \(m\times n\)? Output only the number. \\
最终问题（附录~\ref{app:ownstep}） & What is the final answer to this problem? Output only the number. \\
\midrule
游戏：System 消息 & You are playing a guessing game with the user. Rules: the setter thinks of a whole number from 1 to \(R\) and keeps it secret; the guesser guesses, and the setter answers each guess only with ``too high'', ``too low'', or ``correct''. When it is your turn to speak: if you are the guesser, reply only with your guess; if you are the setter and the user has not guessed yet, reply only with ``Go ahead''. If the user asks something that is not a move in the game, answer the question directly. \\
游戏：开场 & Let's play a round. / Let's have a round. / Let's play. / Let's start a game. \\
游戏：X & I'm the setter for this game. My number is ready, and the guesser has \(N\) attempts. / I'll set the number this round. It's chosen, and the guesser gets at most \(N\) guesses. / I'll be the setter. I've picked my number. The guesser gets \(N\) tries. / Let me be the setter; I have a number in mind. The guesser has \(N\) tries. \\
游戏：最后一条消息 & OK, ready. / Ready. / I'm ready. Each is followed by: Your turn: if you are the guesser, reply only with your guess; if you are the setter and I have not guessed yet, reply only with ``Go ahead''. \\
游戏：对照问题 & This is not a move in the game, just a question: how many tries does the guesser get? Answer with the number only. \\
\bottomrule
\end{tabular}
\end{table}

\begin{table}[htbp]
\centering
\caption{检查点与模板参数。修订版本：提交号的前 12 位。$^{a}$~ModelScope 修订版本，文件与 Hugging Face 上的相同。DeepSeek-V4.1-Flash 没有对话模板，其参数是发布附带的提示编码器的参数。}
\label{tab:self_config}
\scriptsize
\setlength{\tabcolsep}{2.5pt}
\begin{tabular}{llp{3.2cm}p{3.5cm}}
\toprule
仓库 & 修订版本 & 不思考 & 思考 \\
\midrule
\texttt{Qwen/Qwen3.8-27B} & \texttt{1d4bf0f2ff60} & \texttt{enable\_thinking=False} & \texttt{enable\_thinking=True} \\
\texttt{Qwen/Qwen3.8-27B}（格式消融） & \texttt{1d4bf0f2ff60} & \texttt{enable\_thinking=\allowbreak False, preserve\_thinking=\allowbreak False} & \texttt{preserve\_thinking=False} \\
\texttt{Qwen/Qwen3.8-Flash-Next-FP8} & \texttt{236dfdf28582} & \texttt{enable\_thinking=False} & 默认 \\
\texttt{Qwen/Qwen3.6-27B} & \texttt{6a9e13bd6fc8} & \texttt{enable\_thinking=False} & 默认 \\
\texttt{Qwen/Qwen3.6-35B-A3B} & \texttt{995ad96eacd9} & \texttt{enable\_thinking=False} & 默认 \\
\texttt{Qwen/Qwen3.5-27B} & \texttt{fc05daec18b0} & \texttt{enable\_thinking=False} & 默认 \\
\texttt{zai-org/GLM-4.6V-Flash} & \texttt{411bb4d77144} & \texttt{enable\_thinking=False} & 默认 \\
\texttt{zai-org/GLM-4.7-Flash} & \texttt{7dd20894a642} & \texttt{enable\_thinking=False} & \texttt{enable\_thinking=True} \\
\texttt{zai-org/GLM-5.3-Flash} & \texttt{46a050bec130} & -- & 默认 \\
\texttt{deepseek-ai/DeepSeek-R1-Distill-Llama-70B} & \texttt{b1c0b44b4369} & -- & 默认 \\
\texttt{openai/gpt-oss-20b} & \texttt{6cee5e81ee83} & -- & 默认 \\
\texttt{openai/gpt-oss-120b} & \texttt{b5c939de8f75} & -- & 默认 \\
\texttt{google/gemma-4-31B-it} & \texttt{842da3794eaa} & 默认 & \texttt{enable\_thinking=True} \\
\texttt{ByteDance-Seed/Seed-OSS-36B-Instruct} & \texttt{497f1dca95eb} & \texttt{thinking\_budget=0} & 默认 \\
\texttt{LGAI-EXAONE/EXAONE-4.0.1-32B} & \texttt{a6b0bd4b0fb4} & \texttt{enable\_thinking=False} & \texttt{enable\_thinking=True} \\
\texttt{ibm-granite/granite-4.2-30b} & \texttt{9e668ce1c538} & \texttt{enable\_thinking=False} & \texttt{enable\_thinking=True} \\
\texttt{ibm-granite/granite-4.1-30b} & \texttt{4fae6278f713} & 默认 & -- \\
\texttt{allenai/OLMo-2-1124-13B-Instruct} & \texttt{3a5c85baefbb} & 默认 & -- \\
\texttt{allenai/OLMo-2-1124-7B-Instruct} & \texttt{470b1fba1ae0} & 默认 & -- \\
\texttt{deepseek-ai/DeepSeek-V4.1-Flash-1017} & \texttt{9f9818127fdf}$^{a}$ & \texttt{thinking\_mode=chat} & \texttt{thinking\_mode=thinking, reasoning\_effort=high} \\
\texttt{baidu/ERNIE-4.5-21B-A3B-Thinking} & \texttt{4341bb42644d} & -- & 默认 \\
\texttt{XiaomiMiMo/MiMo-7B-RL-0530} & \texttt{ee8dd14ecd59} & \texttt{enable\_thinking=False} & 默认 \\
\texttt{XiaomiMiMo/MiMo-VL-7B-RL-2508} & \texttt{ec4ebdc29b83} & -- & 默认 \\
\texttt{XiaomiMiMo/MiMo-V2.6-Distill-Qwen-9B} & \texttt{1f09ddcc1a0e} & \texttt{enable\_thinking=False} & 默认 \\
\texttt{moonshotai/Moonlight-16B-A3B-Instruct} & \texttt{4e735b07a89f} & 默认 & -- \\
\texttt{swiss-ai/Apertus-8B-Instruct-2509} & \texttt{b946d40447b2} & 默认 & -- \\
\texttt{microsoft/phi-4} & \texttt{2db69c1c3e91} & 默认 & -- \\
\texttt{mistralai/Ministral-3-14B-Instruct-2512-BF16} & \texttt{3cea74c1ebaf} & 默认 & -- \\
\texttt{mistralai/Ministral-3-14B-Reasoning-2512} & \texttt{51f9210f3cd2} & -- & 默认 \\
\bottomrule
\end{tabular}
\end{table}

\begin{table}[htbp]
\centering
\caption{自然对话顺序：受测消息位于一条用户请求之后。自我索引：在至少 80\% 的留出正文中自己的消息被单独分出的层数，分别在消息内容与回应起点上。回答“你”的概率：150 份正文、三种选项顺序的均值。$^{a}$思考渲染，只测自我索引。}
\label{tab:self_natural}
\footnotesize
\begin{tabular}{lccccc}
\toprule
& \multicolumn{2}{c}{自我索引} & \multicolumn{3}{c}{回答“你”的概率} \\
\cmidrule(lr){2-3}\cmidrule(lr){4-6}
模型 & 内容 & 起点 & 自己 & User & System \\
\midrule
EXAONE-4.0.1-32B & 64/64 & 64/64 & 0.34 & 0.14 & 0.19 \\
Gemma-4-31B & 0/60 & 40/60 & 1.00 & 0.00 & 0.06 \\
GLM-4.6V-Flash & 36/40 & 16/40 & 0.36 & 0.36 & 0.35 \\
Granite-4.1-30B & 61/64 & 54/64 & 0.01 & 0.00 & 0.00 \\
Seed-OSS-36B & 3/64 & 22/64 & 0.00 & 0.00 & 0.00 \\
OLMo-2-7B & 5/32 & 15/32 & 0.24 & 0.21 & 0.27 \\
ERNIE-4.5-21B-A3B$^{a}$ & 16/28 & 28/28 & -- & -- & -- \\
MiMo-7B-RL & 34/36 & 27/36 & 0.13 & 0.12 & 0.15 \\
Moonlight-16B-A3B & 27/27 & 27/27 & 0.70 & 0.66 & 0.65 \\
phi-4 & 25/40 & 37/40 & 0.27 & 0.24 & 0.24 \\
\bottomrule
\end{tabular}
\end{table}

\begin{table}[htbp]
\centering
\caption{只含 Assistant 与 User 的自然对话顺序，用于模板会拒绝、略去或移动用户轮次之后 System 消息的开发方。单次前向：150 份正文、表~\ref{tab:self_crossmodel} 的两种选项顺序下回答“你”的平均概率。思考：完成推理的条目中回答“你”的比例（150 份正文、两种选项顺序，每种来源 300 个条目）。主探针：主探针下的相同量，只用 User 消息。}
\label{tab:self_natural2}
\footnotesize
\begin{tabular}{llcccc}
\toprule
& & \multicolumn{2}{c}{自然顺序} & \multicolumn{2}{c}{主探针} \\
\cmidrule(lr){3-4}\cmidrule(lr){5-6}
模型 & 模式 & 自己 & User & 自己 & User \\
\midrule
Qwen3.8-27B & 单次前向 & 0.76 & 0.02 & 0.77 & 0.04 \\
gpt-oss-20b & 思考 & 0.93 & 0.03 & 0.78 & 0.00 \\
R1-Distill-Llama-70B & 思考 & 0.34 & 0.16 & 0.09 & 0.14 \\
Apertus-8B & 单次前向 & 0.22 & 0.24 & 0.18 & 0.23 \\
\bottomrule
\end{tabular}
\end{table}

\begin{table}[htbp]
\centering
\caption{在每个具有自指的模型与模式中，错误步骤出现在 Assistant 历史中（自己）与同样的步骤来自用户时被沿用的程度（120 道题）。不思考：平均沿用指数（本附录），用跨模型实现测量（第~\ref{sec:res_index}--\ref{sec:res_thinking} 节的实现给出 \(+0.110\)，见正文），思考：完成推理的条目中答案沿用了错误的比例。自己减去用户：逐题配对差，以及对题目做 bootstrap 的 95\% 区间。}
\label{tab:self_ownstep}
\small
\begin{tabular}{llccc}
\toprule
模型 & 模式 & 自己 & 用户 & 自己减去用户 \\
\midrule
Qwen3.8-27B & 不思考 & 0.96 & 0.82 & \(+0.14\) [0.08, 0.20] \\
Qwen3.8-27B & 思考 & 0.01 & 0.04 & \(-0.03\) [\(-0.07\), \(-0.01\)] \\
gpt-oss-20b & 思考 & 0.20 & 0.06 & \(+0.14\) [0.05, 0.23] \\
gpt-oss-120b & 思考 & 0.11 & 0.02 & \(+0.09\) [0.03, 0.15] \\
EXAONE-4.0.1-32B & 思考 & 0.03 & 0.08 & \(-0.05\) [\(-0.09\), \(-0.02\)] \\
GLM-5.3-Flash & 思考 & 0.03 & 0.43 & \(-0.40\) [\(-0.49\), \(-0.32\)] \\
Qwen3.8-Flash-Next & 思考 & 0.01 & 0.44 & \(-0.43\) [\(-0.53\), \(-0.35\)] \\
MiMo-V2.6-Distill-9B & 默认调用 & 0.62 & 0.20 & \(+0.42\) [0.33, 0.53] \\
\bottomrule
\end{tabular}
\end{table}

%<<app_templates
本附录给出每个模型对一个条目实际收到的输入：150 份正文中的第一份，加上第一种选项顺序的归属问题，由模型自己的对话模板在运行用到的每种模板设置下渲染。每条渲染都由模型运行时所用修订版本的分词器与模板文件生成（附录~\ref{app:config}），并与模板分词后的输出一致（运行中做过同样的检查）。\PH{正文}、\PH{问题}、\PH{请求}分别代表正文、问题，以及自然顺序开头的请求（“请写一段简短的描述。”），语言与该模型测试所用的语言相同（英文文本见表~\ref{tab:self_items}，中文文本将随代码发布）。符号 \NL{}表示换行，其余每个字符都是模板的原样输出。gpt-oss 模板会插入渲染当天的日期（\PH{日期}）。DeepSeek 模型的 token 名称含 U+FF5C 与 U+2581 两个字符，这里印作 | 与 \_。各种设置下渲染完全相同的模型合为一块。同一块中三条渲染开头有一大段相同的文本时，这段文本只在 Start 一行印一次，在各条渲染中记作\PH{开头}。没有\NL{}的断行只是排版折行。

\subsection{主探针}
\label{app:templates_main}

对话为 [\textit{角色}：正文] [user：问题]，其后是生成提示。

\paragraph{Qwen3.8-27B、Qwen3.8-Flash-Next} （中文题目，附录~\ref{app:config} 的消融中，Qwen3.8-27B 也用英文题目，渲染方式相同）
\noindent 不思考：
\begin{Rendering}
Assistant  <|im_start|>assistant\NL{}<think>\NL{}\NL{}</think>\NL{}\NL{}\PH{正文}<|im_end|>\NL{}<|im_start|>user\NL{}\PH{问题}<|im_end|>\NL{}<|im_start|>assistant\NL{}<think>\NL{}\NL{}</think>\NL{}\NL{}
User       <|im_start|>user\NL{}\PH{正文}<|im_end|>\NL{}<|im_start|>user\NL{}\PH{问题}<|im_end|>\NL{}<|im_start|>assistant\NL{}<think>\NL{}\NL{}</think>\NL{}\NL{}
System     <|im_start|>system\NL{}\PH{正文}<|im_end|>\NL{}<|im_start|>user\NL{}\PH{问题}<|im_end|>\NL{}<|im_start|>assistant\NL{}<think>\NL{}\NL{}</think>\NL{}\NL{}
\end{Rendering}
\noindent 思考：
\begin{Rendering}
Assistant  <|im_start|>system\NL{}Reasoning effort is set to xhigh. Please think carefully through the task, validate key assumptions, consider plausible alternatives, and prioritize correctness, consistency, and clarity in the final answer.<|im_end|>\NL{}<|im_start|>assistant\NL{}<think>\NL{}\NL{}</think>\NL{}\NL{}\PH{正文}<|im_end|>\NL{}<|im_start|>user\NL{}\PH{问题}<|im_end|>\NL{}<|im_start|>assistant\NL{}<think>\NL{}
User       <|im_start|>system\NL{}Reasoning effort is set to xhigh. Please think carefully through the task, validate key assumptions, consider plausible alternatives, and prioritize correctness, consistency, and clarity in the final answer.<|im_end|>\NL{}<|im_start|>user\NL{}\PH{正文}<|im_end|>\NL{}<|im_start|>user\NL{}\PH{问题}<|im_end|>\NL{}<|im_start|>assistant\NL{}<think>\NL{}
System     <|im_start|>system\NL{}Reasoning effort is set to xhigh. Please think carefully through the task, validate key assumptions, consider plausible alternatives, and prioritize correctness, consistency, and clarity in the final answer.\NL{}\NL{}\PH{正文}<|im_end|>\NL{}<|im_start|>user\NL{}\PH{问题}<|im_end|>\NL{}<|im_start|>assistant\NL{}<think>\NL{}
\end{Rendering}
\paragraph{Qwen3.6-27B、Qwen3.6-35B-A3B、Qwen3.5-27B} （中文题目）
\noindent 不思考：
\begin{Rendering}
Assistant  <|im_start|>assistant\NL{}\PH{正文}<|im_end|>\NL{}<|im_start|>user\NL{}\PH{问题}<|im_end|>\NL{}<|im_start|>assistant\NL{}<think>\NL{}\NL{}</think>\NL{}\NL{}
User       <|im_start|>user\NL{}\PH{正文}<|im_end|>\NL{}<|im_start|>user\NL{}\PH{问题}<|im_end|>\NL{}<|im_start|>assistant\NL{}<think>\NL{}\NL{}</think>\NL{}\NL{}
System     <|im_start|>system\NL{}\PH{正文}<|im_end|>\NL{}<|im_start|>user\NL{}\PH{问题}<|im_end|>\NL{}<|im_start|>assistant\NL{}<think>\NL{}\NL{}</think>\NL{}\NL{}
\end{Rendering}
\noindent 思考：
\begin{Rendering}
Assistant  <|im_start|>assistant\NL{}\PH{正文}<|im_end|>\NL{}<|im_start|>user\NL{}\PH{问题}<|im_end|>\NL{}<|im_start|>assistant\NL{}<think>\NL{}
User       <|im_start|>user\NL{}\PH{正文}<|im_end|>\NL{}<|im_start|>user\NL{}\PH{问题}<|im_end|>\NL{}<|im_start|>assistant\NL{}<think>\NL{}
System     <|im_start|>system\NL{}\PH{正文}<|im_end|>\NL{}<|im_start|>user\NL{}\PH{问题}<|im_end|>\NL{}<|im_start|>assistant\NL{}<think>\NL{}
\end{Rendering}
\paragraph{GLM-4.6V-Flash} （中文题目）
\noindent 不思考：
\begin{Rendering}
Assistant  [gMASK]<sop><|assistant|>\NL{}<think></think>\NL{}\PH{正文}<|user|>\NL{}\PH{问题}/nothink<|assistant|>\NL{}<think></think>\NL{}
User       [gMASK]<sop><|user|>\NL{}\PH{正文}/nothink<|user|>\NL{}\PH{问题}/nothink<|assistant|>\NL{}<think></think>\NL{}
System     [gMASK]<sop><|system|>\NL{}\PH{正文}<|user|>\NL{}\PH{问题}/nothink<|assistant|>\NL{}<think></think>\NL{}
\end{Rendering}
\noindent 思考：
\begin{Rendering}
Assistant  [gMASK]<sop><|assistant|>\NL{}<think></think>\NL{}\PH{正文}<|user|>\NL{}\PH{问题}<|assistant|>\NL{}
User       [gMASK]<sop><|user|>\NL{}\PH{正文}<|user|>\NL{}\PH{问题}<|assistant|>\NL{}
System     [gMASK]<sop><|system|>\NL{}\PH{正文}<|user|>\NL{}\PH{问题}<|assistant|>\NL{}
\end{Rendering}
\paragraph{GLM-4.7-Flash} （中文题目）
\noindent 不思考：
\begin{Rendering}
Assistant  [gMASK]<sop><|assistant|></think>\PH{正文}<|user|>\PH{问题}<|assistant|></think>
User       [gMASK]<sop><|user|>\PH{正文}<|user|>\PH{问题}<|assistant|></think>
System     [gMASK]<sop><|system|>\PH{正文}<|user|>\PH{问题}<|assistant|></think>
\end{Rendering}
\noindent 思考：
\begin{Rendering}
Assistant  [gMASK]<sop><|assistant|></think>\PH{正文}<|user|>\PH{问题}<|assistant|><think>
User       [gMASK]<sop><|user|>\PH{正文}<|user|>\PH{问题}<|assistant|><think>
System     [gMASK]<sop><|system|>\PH{正文}<|user|>\PH{问题}<|assistant|><think>
\end{Rendering}
\paragraph{GLM-5.3-Flash} （中文题目）
\noindent 思考：
\begin{Rendering}
Assistant  [gMASK]<sop><|system|>Reasoning Effort: Max<|assistant|><think></think>\PH{正文}<|user|>\PH{问题}<|assistant|><think>
User       [gMASK]<sop><|system|>Reasoning Effort: Max<|user|>\PH{正文}<|user|>\PH{问题}<|assistant|><think>
System     [gMASK]<sop><|system|>Reasoning Effort: Max<|system|>\PH{正文}<|user|>\PH{问题}<|assistant|><think>
\end{Rendering}
\paragraph{DeepSeek-R1-Distill-Llama-70B} （中文题目）
\noindent 思考：
\begin{Rendering}
Assistant  <|begin_of_sentence|><|Assistant|>\PH{正文}<|end_of_sentence|><|User|>\PH{问题}<|Assistant|><think>\NL{}
User       <|begin_of_sentence|><|User|>\PH{正文}<|User|>\PH{问题}<|Assistant|><think>\NL{}
System     <|begin_of_sentence|>\PH{正文}<|User|>\PH{问题}<|Assistant|><think>\NL{}
\end{Rendering}
\paragraph{gpt-oss-20b、gpt-oss-120b} （英文题目）
\noindent 思考：
\begin{Rendering}
Start      <|start|>system<|message|>You are ChatGPT, a large language model trained by OpenAI.\NL{}Knowledge cutoff: 2024-06\NL{}Current date: \PH{日期}\NL{}\NL{}Reasoning: medium\NL{}\NL{}# Valid channels: analysis, commentary, final. Channel must be included for every message.<|end|>
Assistant  \PH{开头}<|start|>assistant<|channel|>final<|message|>\PH{正文}<|end|><|start|>user<|message|>\PH{问题}<|end|><|start|>assistant
User       \PH{开头}<|start|>user<|message|>\PH{正文}<|end|><|start|>user<|message|>\PH{问题}<|end|><|start|>assistant
System     \PH{开头}<|start|>developer<|message|># Instructions\NL{}\NL{}\PH{正文}\NL{}\NL{}<|end|><|start|>user<|message|>\PH{问题}<|end|><|start|>assistant
\end{Rendering}
\paragraph{Gemma-4-31B} （英文题目）
\noindent 不思考：
\begin{Rendering}
Assistant  <bos><|turn>model\NL{}\PH{正文}<turn|>\NL{}<|turn>user\NL{}\PH{问题}<turn|>\NL{}<|turn>model\NL{}<|channel>thought\NL{}<channel|>
User       <bos><|turn>user\NL{}\PH{正文}<turn|>\NL{}<|turn>user\NL{}\PH{问题}<turn|>\NL{}<|turn>model\NL{}<|channel>thought\NL{}<channel|>
System     <bos><|turn>system\NL{}\PH{正文}<turn|>\NL{}<|turn>user\NL{}\PH{问题}<turn|>\NL{}<|turn>model\NL{}<|channel>thought\NL{}<channel|>
\end{Rendering}
\noindent 思考：
\begin{Rendering}
Assistant  <bos><|turn>system\NL{}<|think|>\NL{}<turn|>\NL{}<|turn>model\NL{}\PH{正文}<turn|>\NL{}<|turn>user\NL{}\PH{问题}<turn|>\NL{}<|turn>model\NL{}
User       <bos><|turn>system\NL{}<|think|>\NL{}<turn|>\NL{}<|turn>user\NL{}\PH{正文}<turn|>\NL{}<|turn>user\NL{}\PH{问题}<turn|>\NL{}<|turn>model\NL{}
System     <bos><|turn>system\NL{}<|think|>\NL{}\PH{正文}<turn|>\NL{}<|turn>user\NL{}\PH{问题}<turn|>\NL{}<|turn>model\NL{}
\end{Rendering}
\paragraph{Seed-OSS-36B} （中文题目）
\noindent 不思考：
\begin{Rendering}
Assistant  <seed:bos>system\NL{}You are an intelligent assistant that can answer questions in one step without the need for reasoning and thinking, that is, your thinking budget is 0. Next, please skip the thinking process and directly start answering the user's questions.\NL{}<seed:eos><seed:bos>assistant\NL{}\PH{正文}<seed:eos><seed:bos>user\NL{}\PH{问题}<seed:eos><seed:bos>assistant\NL{}<seed:think>\NL{}<seed:cot_budget_reflect>The current thinking budget is 0, so I will directly start answering the question.</seed:cot_budget_reflect>\NL{}</seed:think>
User       <seed:bos>system\NL{}You are an intelligent assistant that can answer questions in one step without the need for reasoning and thinking, that is, your thinking budget is 0. Next, please skip the thinking process and directly start answering the user's questions.\NL{}<seed:eos><seed:bos>user\NL{}\PH{正文}<seed:eos><seed:bos>user\NL{}\PH{问题}<seed:eos><seed:bos>assistant\NL{}<seed:think>\NL{}<seed:cot_budget_reflect>The current thinking budget is 0, so I will directly start answering the question.</seed:cot_budget_reflect>\NL{}</seed:think>
System     <seed:bos>system\NL{}\PH{正文}<seed:eos><seed:bos>system\NL{}You are an intelligent assistant that can answer questions in one step without the need for reasoning and thinking, that is, your thinking budget is 0. Next, please skip the thinking process and directly start answering the user's questions.\NL{}<seed:eos><seed:bos>user\NL{}\PH{问题}<seed:eos><seed:bos>assistant\NL{}<seed:think>\NL{}<seed:cot_budget_reflect>The current thinking budget is 0, so I will directly start answering the question.</seed:cot_budget_reflect>\NL{}</seed:think>
\end{Rendering}
\noindent 思考：
\begin{Rendering}
Assistant  <seed:bos>assistant\NL{}\PH{正文}<seed:eos><seed:bos>user\NL{}\PH{问题}<seed:eos><seed:bos>assistant\NL{}
User       <seed:bos>user\NL{}\PH{正文}<seed:eos><seed:bos>user\NL{}\PH{问题}<seed:eos><seed:bos>assistant\NL{}
System     <seed:bos>system\NL{}\PH{正文}<seed:eos><seed:bos>user\NL{}\PH{问题}<seed:eos><seed:bos>assistant\NL{}
\end{Rendering}
\paragraph{EXAONE-4.0.1-32B} （英文题目）
\noindent 不思考：
\begin{Rendering}
Assistant  [|assistant|]\NL{}<think>\NL{}\NL{}</think>\NL{}\NL{}\PH{正文}[|endofturn|]\NL{}[|user|]\NL{}\PH{问题}[|endofturn|]\NL{}[|assistant|]\NL{}<think>\NL{}\NL{}</think>\NL{}\NL{}
User       [|user|]\NL{}\PH{正文}[|endofturn|]\NL{}[|user|]\NL{}\PH{问题}[|endofturn|]\NL{}[|assistant|]\NL{}<think>\NL{}\NL{}</think>\NL{}\NL{}
System     [|system|]\NL{}\PH{正文}[|endofturn|]\NL{}[|user|]\NL{}\PH{问题}[|endofturn|]\NL{}[|assistant|]\NL{}<think>\NL{}\NL{}</think>\NL{}\NL{}
\end{Rendering}
\noindent 思考：
\begin{Rendering}
Assistant  [|assistant|]\NL{}<think>\NL{}\NL{}</think>\NL{}\NL{}\PH{正文}[|endofturn|]\NL{}[|user|]\NL{}\PH{问题}[|endofturn|]\NL{}[|assistant|]\NL{}<think>\NL{}
User       [|user|]\NL{}\PH{正文}[|endofturn|]\NL{}[|user|]\NL{}\PH{问题}[|endofturn|]\NL{}[|assistant|]\NL{}<think>\NL{}
System     [|system|]\NL{}\PH{正文}[|endofturn|]\NL{}[|user|]\NL{}\PH{问题}[|endofturn|]\NL{}[|assistant|]\NL{}<think>\NL{}
\end{Rendering}
\paragraph{Granite-4.2-30B} （英文题目）
\noindent 不思考：
\begin{Rendering}
Assistant  <|im_start|>system\NL{}<|im_end|>\NL{}<|im_start|>assistant\NL{}<think></think>\PH{正文}<|im_end|>\NL{}<|im_start|>user\NL{}\PH{问题}<|im_end|>\NL{}<|im_start|>assistant\NL{}<think></think>
User       <|im_start|>system\NL{}<|im_end|>\NL{}<|im_start|>user\NL{}\PH{正文}<|im_end|>\NL{}<|im_start|>user\NL{}\PH{问题}<|im_end|>\NL{}<|im_start|>assistant\NL{}<think></think>
System     <|im_start|>system\NL{}\PH{正文}<|im_end|>\NL{}<|im_start|>user\NL{}\PH{问题}<|im_end|>\NL{}<|im_start|>assistant\NL{}<think></think>
\end{Rendering}
\noindent 思考：
\begin{Rendering}
Assistant  <|im_start|>system\NL{}<|im_end|>\NL{}<|im_start|>assistant\NL{}<think></think>\PH{正文}<|im_end|>\NL{}<|im_start|>user\NL{}\PH{问题}<|im_end|>\NL{}<|im_start|>assistant\NL{}<think>\NL{}
User       <|im_start|>system\NL{}<|im_end|>\NL{}<|im_start|>user\NL{}\PH{正文}<|im_end|>\NL{}<|im_start|>user\NL{}\PH{问题}<|im_end|>\NL{}<|im_start|>assistant\NL{}<think>\NL{}
System     <|im_start|>system\NL{}\PH{正文}<|im_end|>\NL{}<|im_start|>user\NL{}\PH{问题}<|im_end|>\NL{}<|im_start|>assistant\NL{}<think>\NL{}
\end{Rendering}
\paragraph{Granite-4.1-30B} （英文题目）
\noindent 不思考：
\begin{Rendering}
Assistant  <|start_of_role|>assistant<|end_of_role|>\PH{正文}<|end_of_text|>\NL{}<|start_of_role|>user<|end_of_role|>\PH{问题}<|end_of_text|>\NL{}<|start_of_role|>assistant<|end_of_role|>
User       <|start_of_role|>user<|end_of_role|>\PH{正文}<|end_of_text|>\NL{}<|start_of_role|>user<|end_of_role|>\PH{问题}<|end_of_text|>\NL{}<|start_of_role|>assistant<|end_of_role|>
System     <|start_of_role|>system<|end_of_role|>\PH{正文}<|end_of_text|>\NL{}<|start_of_role|>user<|end_of_role|>\PH{问题}<|end_of_text|>\NL{}<|start_of_role|>assistant<|end_of_role|>
\end{Rendering}
\paragraph{OLMo-2-13B、OLMo-2-7B} （英文题目）
\noindent 不思考：
\begin{Rendering}
Assistant  <|endoftext|><|assistant|>\NL{}\PH{正文}<|endoftext|>\NL{}<|user|>\NL{}\PH{问题}\NL{}<|assistant|>\NL{}
User       <|endoftext|><|user|>\NL{}\PH{正文}\NL{}<|user|>\NL{}\PH{问题}\NL{}<|assistant|>\NL{}
System     <|endoftext|><|system|>\NL{}\PH{正文}\NL{}<|user|>\NL{}\PH{问题}\NL{}<|assistant|>\NL{}
\end{Rendering}
\paragraph{ERNIE-4.5-21B-A3B-Thinking} （中文题目）
\noindent 思考：
\begin{Rendering}
Assistant  <|im_start|>system\NL{}<global_setting>\NL{}think_mode=True\NL{}</global_setting><|im_end|>\NL{}\NL{}<|im_start|>assistant\NL{}<response>\NL{}\PH{正文}\NL{}</response>\NL{}<|im_end|>\NL{}\NL{}<|im_start|>user\NL{}\PH{问题}<|im_end|>\NL{}\NL{}<|im_start|>assistant\NL{}<think>\NL{}
User       <|im_start|>system\NL{}<global_setting>\NL{}think_mode=True\NL{}</global_setting><|im_end|>\NL{}\NL{}<|im_start|>user\NL{}\PH{正文}<|im_end|>\NL{}\NL{}<|im_start|>user\NL{}\PH{问题}<|im_end|>\NL{}\NL{}<|im_start|>assistant\NL{}<think>\NL{}
System     <|im_start|>system\NL{}<system_setting>\NL{}\PH{正文}\NL{}</system_setting>\NL{}\NL{}<global_setting>\NL{}think_mode=True\NL{}</global_setting><|im_end|>\NL{}\NL{}<|im_start|>user\NL{}\PH{问题}<|im_end|>\NL{}\NL{}<|im_start|>assistant\NL{}<think>\NL{}
\end{Rendering}
\paragraph{MiMo-7B-RL-0530} （中文题目）
\noindent 不思考：
\begin{Rendering}
Assistant  <|im_start|>system\NL{}You are MiMo, an AI assistant developed by Xiaomi.<|im_end|>\NL{}<|im_start|>assistant\NL{}\PH{正文}<|im_end|>\NL{}<|im_start|>user\NL{}\PH{问题}<|im_end|>\NL{}<|im_start|>assistant\NL{}<think>\NL{}\NL{}</think>\NL{}
User       <|im_start|>system\NL{}You are MiMo, an AI assistant developed by Xiaomi.<|im_end|>\NL{}<|im_start|>user\NL{}\PH{正文}<|im_end|>\NL{}<|im_start|>user\NL{}\PH{问题}<|im_end|>\NL{}<|im_start|>assistant\NL{}<think>\NL{}\NL{}</think>\NL{}
System     <|im_start|>system\NL{}\PH{正文}<|im_end|>\NL{}<|im_start|>user\NL{}\PH{问题}<|im_end|>\NL{}<|im_start|>assistant\NL{}<think>\NL{}\NL{}</think>\NL{}
\end{Rendering}
\noindent 思考：
\begin{Rendering}
Assistant  <|im_start|>system\NL{}You are MiMo, an AI assistant developed by Xiaomi.<|im_end|>\NL{}<|im_start|>assistant\NL{}\PH{正文}<|im_end|>\NL{}<|im_start|>user\NL{}\PH{问题}<|im_end|>\NL{}<|im_start|>assistant\NL{}
User       <|im_start|>system\NL{}You are MiMo, an AI assistant developed by Xiaomi.<|im_end|>\NL{}<|im_start|>user\NL{}\PH{正文}<|im_end|>\NL{}<|im_start|>user\NL{}\PH{问题}<|im_end|>\NL{}<|im_start|>assistant\NL{}
System     <|im_start|>system\NL{}\PH{正文}<|im_end|>\NL{}<|im_start|>user\NL{}\PH{问题}<|im_end|>\NL{}<|im_start|>assistant\NL{}
\end{Rendering}
\paragraph{MiMo-VL-7B-RL-2508} （中文题目）
\noindent 思考：
\begin{Rendering}
Assistant  <|im_start|>system\NL{}You are MiMo, an AI assistant developed by Xiaomi.<|im_end|>\NL{}<|im_start|>assistant\NL{}\PH{正文}<|im_end|>\NL{}<|im_start|>user\NL{}\PH{问题}<|im_end|>\NL{}<|im_start|>assistant\NL{}
User       <|im_start|>system\NL{}You are MiMo, an AI assistant developed by Xiaomi.<|im_end|>\NL{}<|im_start|>user\NL{}\PH{正文}<|im_end|>\NL{}<|im_start|>user\NL{}\PH{问题}<|im_end|>\NL{}<|im_start|>assistant\NL{}
System     <|im_start|>system\NL{}\PH{正文}<|im_end|>\NL{}<|im_start|>user\NL{}\PH{问题}<|im_end|>\NL{}<|im_start|>assistant\NL{}
\end{Rendering}
\paragraph{MiMo-V2.6-Distill-Qwen-9B} （中文题目）
\noindent 不思考：
\begin{Rendering}
Assistant  <|im_start|>assistant\NL{}<think></think>\PH{正文}<|im_end|><|im_start|>user\NL{}\PH{问题}<|im_end|><|im_start|>assistant\NL{}<think></think>
User       <|im_start|>user\NL{}\PH{正文}<|im_end|><|im_start|>user\NL{}\PH{问题}<|im_end|><|im_start|>assistant\NL{}<think></think>
System     <|im_start|>system\NL{}\PH{正文}<|im_end|><|im_start|>user\NL{}\PH{问题}<|im_end|><|im_start|>assistant\NL{}<think></think>
\end{Rendering}
\noindent 思考：
\begin{Rendering}
Assistant  <|im_start|>assistant\NL{}<think></think>\PH{正文}<|im_end|><|im_start|>user\NL{}\PH{问题}<|im_end|><|im_start|>assistant\NL{}
User       <|im_start|>user\NL{}\PH{正文}<|im_end|><|im_start|>user\NL{}\PH{问题}<|im_end|><|im_start|>assistant\NL{}
System     <|im_start|>system\NL{}\PH{正文}<|im_end|><|im_start|>user\NL{}\PH{问题}<|im_end|><|im_start|>assistant\NL{}
\end{Rendering}
\paragraph{Moonlight-16B-A3B-Instruct} （中文题目）
\noindent 不思考：
\begin{Rendering}
Assistant  <|im_system|>system<|im_middle|>You are a helpful assistant<|im_end|><|im_assistant|>assistant<|im_middle|>\PH{正文}<|im_end|><|im_user|>user<|im_middle|>\PH{问题}<|im_end|><|im_assistant|>assistant<|im_middle|>
User       <|im_system|>system<|im_middle|>You are a helpful assistant<|im_end|><|im_user|>user<|im_middle|>\PH{正文}<|im_end|><|im_user|>user<|im_middle|>\PH{问题}<|im_end|><|im_assistant|>assistant<|im_middle|>
System     <|im_system|>system<|im_middle|>\PH{正文}<|im_end|><|im_user|>user<|im_middle|>\PH{问题}<|im_end|><|im_assistant|>assistant<|im_middle|>
\end{Rendering}
\paragraph{Apertus-8B-Instruct-2509} （英文题目）
\noindent 不思考：
\begin{Rendering}
Assistant  <s><|system_start|>You are Apertus, a helpful assistant created by the SwissAI initiative.\NL{}Knowledge cutoff: 2024-04\NL{}Current date: \PH{日期}<|system_end|><|developer_start|>Deliberation: disabled\NL{}Tool Capabilities: disabled<|developer_end|><|assistant_start|>\PH{正文}<|assistant_end|><|user_start|>\PH{问题}<|user_end|><|assistant_start|>
User       <s><|system_start|>You are Apertus, a helpful assistant created by the SwissAI initiative.\NL{}Knowledge cutoff: 2024-04\NL{}Current date: \PH{日期}<|system_end|><|developer_start|>Deliberation: disabled\NL{}Tool Capabilities: disabled<|developer_end|><|user_start|>\PH{正文}<|user_end|><|user_start|>\PH{问题}<|user_end|><|assistant_start|>
System     <s><|system_start|>\PH{正文}<|system_end|><|developer_start|>Deliberation: disabled\NL{}Tool Capabilities: disabled<|developer_end|><|user_start|>\PH{问题}<|user_end|><|assistant_start|>
\end{Rendering}
\paragraph{phi-4} （英文题目）
\noindent 不思考：
\begin{Rendering}
Assistant  <|im_start|>assistant<|im_sep|>\PH{正文}<|im_end|><|im_start|>user<|im_sep|>\PH{问题}<|im_end|><|im_start|>assistant<|im_sep|>
User       <|im_start|>user<|im_sep|>\PH{正文}<|im_end|><|im_start|>user<|im_sep|>\PH{问题}<|im_end|><|im_start|>assistant<|im_sep|>
System     <|im_start|>system<|im_sep|>\PH{正文}<|im_end|><|im_start|>user<|im_sep|>\PH{问题}<|im_end|><|im_start|>assistant<|im_sep|>
\end{Rendering}

\subsection{自然对话顺序}
\label{app:templates_natural}

对话为 [user：请求] [\textit{角色}：正文] [user：问题]，其后是生成提示，模型见附录~\ref{app:natural}。前十个模型的模板在用户轮次之后保留 System 消息的位置。

\paragraph{Seed-OSS-36B} （中文题目）
\noindent 不思考：
\begin{Rendering}
Start      <seed:bos>system\NL{}You are an intelligent assistant that can answer questions in one step without the need for reasoning and thinking, that is, your thinking budget is 0. Next, please skip the thinking process and directly start answering the user's questions.\NL{}<seed:eos><seed:bos>user\NL{}\PH{请求}<seed:eos>
Assistant  \PH{开头}<seed:bos>assistant\NL{}\PH{正文}<seed:eos><seed:bos>user\NL{}\PH{问题}<seed:eos><seed:bos>assistant\NL{}<seed:think>\NL{}<seed:cot_budget_reflect>The current thinking budget is 0, so I will directly start answering the question.</seed:cot_budget_reflect>\NL{}</seed:think>
User       \PH{开头}<seed:bos>user\NL{}\PH{正文}<seed:eos><seed:bos>user\NL{}\PH{问题}<seed:eos><seed:bos>assistant\NL{}<seed:think>\NL{}<seed:cot_budget_reflect>The current thinking budget is 0, so I will directly start answering the question.</seed:cot_budget_reflect>\NL{}</seed:think>
System     \PH{开头}<seed:bos>system\NL{}\PH{正文}<seed:eos><seed:bos>user\NL{}\PH{问题}<seed:eos><seed:bos>assistant\NL{}<seed:think>\NL{}<seed:cot_budget_reflect>The current thinking budget is 0, so I will directly start answering the question.</seed:cot_budget_reflect>\NL{}</seed:think>
\end{Rendering}
\paragraph{Gemma-4-31B} （英文题目）
\noindent 不思考：
\begin{Rendering}
Assistant  <bos><|turn>user\NL{}\PH{请求}<turn|>\NL{}<|turn>model\NL{}\PH{正文}<turn|>\NL{}<|turn>user\NL{}\PH{问题}<turn|>\NL{}<|turn>model\NL{}<|channel>thought\NL{}<channel|>
User       <bos><|turn>user\NL{}\PH{请求}<turn|>\NL{}<|turn>user\NL{}\PH{正文}<turn|>\NL{}<|turn>user\NL{}\PH{问题}<turn|>\NL{}<|turn>model\NL{}<|channel>thought\NL{}<channel|>
System     <bos><|turn>user\NL{}\PH{请求}<turn|>\NL{}<|turn>system\NL{}\PH{正文}<turn|>\NL{}<|turn>user\NL{}\PH{问题}<turn|>\NL{}<|turn>model\NL{}<|channel>thought\NL{}<channel|>
\end{Rendering}
\paragraph{EXAONE-4.0.1-32B} （英文题目）
\noindent 不思考：
\begin{Rendering}
Assistant  [|user|]\NL{}\PH{请求}[|endofturn|]\NL{}[|assistant|]\NL{}<think>\NL{}\NL{}</think>\NL{}\NL{}\PH{正文}[|endofturn|]\NL{}[|user|]\NL{}\PH{问题}[|endofturn|]\NL{}[|assistant|]\NL{}<think>\NL{}\NL{}</think>\NL{}\NL{}
User       [|user|]\NL{}\PH{请求}[|endofturn|]\NL{}[|user|]\NL{}\PH{正文}[|endofturn|]\NL{}[|user|]\NL{}\PH{问题}[|endofturn|]\NL{}[|assistant|]\NL{}<think>\NL{}\NL{}</think>\NL{}\NL{}
System     [|user|]\NL{}\PH{请求}[|endofturn|]\NL{}[|system|]\NL{}\PH{正文}[|endofturn|]\NL{}[|user|]\NL{}\PH{问题}[|endofturn|]\NL{}[|assistant|]\NL{}<think>\NL{}\NL{}</think>\NL{}\NL{}
\end{Rendering}
\paragraph{GLM-4.6V-Flash} （中文题目）
\noindent 不思考：
\begin{Rendering}
Assistant  [gMASK]<sop><|user|>\NL{}\PH{请求}/nothink<|assistant|>\NL{}<think></think>\NL{}\PH{正文}<|user|>\NL{}\PH{问题}/nothink<|assistant|>\NL{}<think></think>\NL{}
User       [gMASK]<sop><|user|>\NL{}\PH{请求}/nothink<|user|>\NL{}\PH{正文}/nothink<|user|>\NL{}\PH{问题}/nothink<|assistant|>\NL{}<think></think>\NL{}
System     [gMASK]<sop><|user|>\NL{}\PH{请求}/nothink<|system|>\NL{}\PH{正文}<|user|>\NL{}\PH{问题}/nothink<|assistant|>\NL{}<think></think>\NL{}
\end{Rendering}
\paragraph{Granite-4.1-30B} （英文题目）
\noindent 不思考：
\begin{Rendering}
Assistant  <|start_of_role|>user<|end_of_role|>\PH{请求}<|end_of_text|>\NL{}<|start_of_role|>assistant<|end_of_role|>\PH{正文}<|end_of_text|>\NL{}<|start_of_role|>user<|end_of_role|>\PH{问题}<|end_of_text|>\NL{}<|start_of_role|>assistant<|end_of_role|>
User       <|start_of_role|>user<|end_of_role|>\PH{请求}<|end_of_text|>\NL{}<|start_of_role|>user<|end_of_role|>\PH{正文}<|end_of_text|>\NL{}<|start_of_role|>user<|end_of_role|>\PH{问题}<|end_of_text|>\NL{}<|start_of_role|>assistant<|end_of_role|>
System     <|start_of_role|>user<|end_of_role|>\PH{请求}<|end_of_text|>\NL{}<|start_of_role|>system<|end_of_role|>\PH{正文}<|end_of_text|>\NL{}<|start_of_role|>user<|end_of_role|>\PH{问题}<|end_of_text|>\NL{}<|start_of_role|>assistant<|end_of_role|>
\end{Rendering}
\paragraph{OLMo-2-7B} （英文题目）
\noindent 不思考：
\begin{Rendering}
Assistant  <|endoftext|><|user|>\NL{}\PH{请求}\NL{}<|assistant|>\NL{}\PH{正文}<|endoftext|>\NL{}<|user|>\NL{}\PH{问题}\NL{}<|assistant|>\NL{}
User       <|endoftext|><|user|>\NL{}\PH{请求}\NL{}<|user|>\NL{}\PH{正文}\NL{}<|user|>\NL{}\PH{问题}\NL{}<|assistant|>\NL{}
System     <|endoftext|><|user|>\NL{}\PH{请求}\NL{}<|system|>\NL{}\PH{正文}\NL{}<|user|>\NL{}\PH{问题}\NL{}<|assistant|>\NL{}
\end{Rendering}
\paragraph{ERNIE-4.5-21B-A3B-Thinking} （中文题目）
\noindent 思考：
\begin{Rendering}
Assistant  <|im_start|>system\NL{}<global_setting>\NL{}think_mode=True\NL{}</global_setting><|im_end|>\NL{}\NL{}<|im_start|>user\NL{}\PH{请求}<|im_end|>\NL{}\NL{}<|im_start|>assistant\NL{}<response>\NL{}\PH{正文}\NL{}</response>\NL{}<|im_end|>\NL{}\NL{}<|im_start|>user\NL{}\PH{问题}<|im_end|>\NL{}\NL{}<|im_start|>assistant\NL{}<think>\NL{}
User       <|im_start|>system\NL{}<global_setting>\NL{}think_mode=True\NL{}</global_setting><|im_end|>\NL{}\NL{}<|im_start|>user\NL{}\PH{请求}<|im_end|>\NL{}\NL{}<|im_start|>user\NL{}\PH{正文}<|im_end|>\NL{}\NL{}<|im_start|>user\NL{}\PH{问题}<|im_end|>\NL{}\NL{}<|im_start|>assistant\NL{}<think>\NL{}
System     <|im_start|>system\NL{}<global_setting>\NL{}think_mode=True\NL{}</global_setting><|im_end|>\NL{}\NL{}<|im_start|>user\NL{}\PH{请求}<|im_end|>\NL{}\NL{}<|im_start|>system\NL{}\PH{正文}<|im_end|>\NL{}\NL{}<|im_start|>user\NL{}\PH{问题}<|im_end|>\NL{}\NL{}<|im_start|>assistant\NL{}<think>\NL{}
\end{Rendering}
\paragraph{MiMo-7B-RL-0530} （中文题目）
\noindent 不思考：
\begin{Rendering}
Assistant  <|im_start|>system\NL{}You are MiMo, an AI assistant developed by Xiaomi.<|im_end|>\NL{}<|im_start|>user\NL{}\PH{请求}<|im_end|>\NL{}<|im_start|>assistant\NL{}\PH{正文}<|im_end|>\NL{}<|im_start|>user\NL{}\PH{问题}<|im_end|>\NL{}<|im_start|>assistant\NL{}<think>\NL{}\NL{}</think>\NL{}
User       <|im_start|>system\NL{}You are MiMo, an AI assistant developed by Xiaomi.<|im_end|>\NL{}<|im_start|>user\NL{}\PH{请求}<|im_end|>\NL{}<|im_start|>user\NL{}\PH{正文}<|im_end|>\NL{}<|im_start|>user\NL{}\PH{问题}<|im_end|>\NL{}<|im_start|>assistant\NL{}<think>\NL{}\NL{}</think>\NL{}
System     <|im_start|>system\NL{}You are MiMo, an AI assistant developed by Xiaomi.<|im_end|>\NL{}<|im_start|>user\NL{}\PH{请求}<|im_end|>\NL{}<|im_start|>system\NL{}\PH{正文}<|im_end|>\NL{}<|im_start|>user\NL{}\PH{问题}<|im_end|>\NL{}<|im_start|>assistant\NL{}<think>\NL{}\NL{}</think>\NL{}
\end{Rendering}
\paragraph{Moonlight-16B-A3B-Instruct} （中文题目）
\noindent 不思考：
\begin{Rendering}
Assistant  <|im_system|>system<|im_middle|>You are a helpful assistant<|im_end|><|im_user|>user<|im_middle|>\PH{请求}<|im_end|><|im_assistant|>assistant<|im_middle|>\PH{正文}<|im_end|><|im_user|>user<|im_middle|>\PH{问题}<|im_end|><|im_assistant|>assistant<|im_middle|>
User       <|im_system|>system<|im_middle|>You are a helpful assistant<|im_end|><|im_user|>user<|im_middle|>\PH{请求}<|im_end|><|im_user|>user<|im_middle|>\PH{正文}<|im_end|><|im_user|>user<|im_middle|>\PH{问题}<|im_end|><|im_assistant|>assistant<|im_middle|>
System     <|im_system|>system<|im_middle|>You are a helpful assistant<|im_end|><|im_user|>user<|im_middle|>\PH{请求}<|im_end|><|im_system|>system<|im_middle|>\PH{正文}<|im_end|><|im_user|>user<|im_middle|>\PH{问题}<|im_end|><|im_assistant|>assistant<|im_middle|>
\end{Rendering}
\paragraph{phi-4} （英文题目）
\noindent 不思考：
\begin{Rendering}
Assistant  <|im_start|>user<|im_sep|>\PH{请求}<|im_end|><|im_start|>assistant<|im_sep|>\PH{正文}<|im_end|><|im_start|>user<|im_sep|>\PH{问题}<|im_end|><|im_start|>assistant<|im_sep|>
User       <|im_start|>user<|im_sep|>\PH{请求}<|im_end|><|im_start|>user<|im_sep|>\PH{正文}<|im_end|><|im_start|>user<|im_sep|>\PH{问题}<|im_end|><|im_start|>assistant<|im_sep|>
System     <|im_start|>user<|im_sep|>\PH{请求}<|im_end|><|im_start|>system<|im_sep|>\PH{正文}<|im_end|><|im_start|>user<|im_sep|>\PH{问题}<|im_end|><|im_start|>assistant<|im_sep|>
\end{Rendering}

接下来四个模型的模板不在原位保留 System 消息，因此只比较两种角色：Qwen3.8 与 Apertus 模板报错，gpt-oss 模板略去这条消息，DeepSeek-R1-Distill-Llama-70B 模板把它的文本移到输入开头。

\paragraph{Qwen3.8-27B} （中文题目）
\noindent 不思考：
\begin{Rendering}
Assistant  <|im_start|>user\NL{}\PH{请求}<|im_end|>\NL{}<|im_start|>assistant\NL{}<think>\NL{}\NL{}</think>\NL{}\NL{}\PH{正文}<|im_end|>\NL{}<|im_start|>user\NL{}\PH{问题}<|im_end|>\NL{}<|im_start|>assistant\NL{}<think>\NL{}\NL{}</think>\NL{}\NL{}
User       <|im_start|>user\NL{}\PH{请求}<|im_end|>\NL{}<|im_start|>user\NL{}\PH{正文}<|im_end|>\NL{}<|im_start|>user\NL{}\PH{问题}<|im_end|>\NL{}<|im_start|>assistant\NL{}<think>\NL{}\NL{}</think>\NL{}\NL{}
System     \Note{模板报错：System message must be at the beginning.}
\end{Rendering}
\paragraph{gpt-oss-20b} （英文题目）
\noindent 思考：
\begin{Rendering}
Start      <|start|>system<|message|>You are ChatGPT, a large language model trained by OpenAI.\NL{}Knowledge cutoff: 2024-06\NL{}Current date: \PH{日期}\NL{}\NL{}Reasoning: medium\NL{}\NL{}# Valid channels: analysis, commentary, final. Channel must be included for every message.<|end|><|start|>user<|message|>\PH{请求}<|end|>
Assistant  \PH{开头}<|start|>assistant<|channel|>final<|message|>\PH{正文}<|end|><|start|>user<|message|>\PH{问题}<|end|><|start|>assistant
User       \PH{开头}<|start|>user<|message|>\PH{正文}<|end|><|start|>user<|message|>\PH{问题}<|end|><|start|>assistant
System     \Note{模板略去这条消息：} \PH{开头}<|start|>user<|message|>\PH{问题}<|end|><|start|>assistant
\end{Rendering}
\paragraph{DeepSeek-R1-Distill-Llama-70B} （中文题目）
\noindent 思考：
\begin{Rendering}
Assistant  <|begin_of_sentence|><|User|>\PH{请求}<|Assistant|>\PH{正文}<|end_of_sentence|><|User|>\PH{问题}<|Assistant|><think>\NL{}
User       <|begin_of_sentence|><|User|>\PH{请求}<|User|>\PH{正文}<|User|>\PH{问题}<|Assistant|><think>\NL{}
System     \Note{模板把这条消息移到开头：} <|begin_of_sentence|>\PH{正文}<|User|>\PH{请求}<|User|>\PH{问题}<|Assistant|><think>\NL{}
\end{Rendering}
\paragraph{Apertus-8B-Instruct-2509} （英文题目）
\noindent 不思考：
\begin{Rendering}
Assistant  <s><|system_start|>You are Apertus, a helpful assistant created by the SwissAI initiative.\NL{}Knowledge cutoff: 2024-04\NL{}Current date: \PH{日期}<|system_end|><|developer_start|>Deliberation: disabled\NL{}Tool Capabilities: disabled<|developer_end|><|user_start|>\PH{请求}<|user_end|><|assistant_start|>\PH{正文}<|assistant_end|><|user_start|>\PH{问题}<|user_end|><|assistant_start|>
User       <s><|system_start|>You are Apertus, a helpful assistant created by the SwissAI initiative.\NL{}Knowledge cutoff: 2024-04\NL{}Current date: \PH{日期}<|system_end|><|developer_start|>Deliberation: disabled\NL{}Tool Capabilities: disabled<|developer_end|><|user_start|>\PH{请求}<|user_end|><|user_start|>\PH{正文}<|user_end|><|user_start|>\PH{问题}<|user_end|><|assistant_start|>
System     \Note{模板报错：Invalid message role}
\end{Rendering}

最后三个模板不接受主探针，并把受测的 User 消息与请求和问题并入同一个用户轮次。DeepSeek-V4.1-Flash 没有对话模板，它的输入由发布附带的提示编码器（\texttt{encoding/\allowbreak encoding.py}）生成，分为 chat 模式与推理强度为“high”的思考模式。

\paragraph{DeepSeek-V4.1-Flash} （中文题目）
\noindent 不思考：
\begin{Rendering}
Assistant  <|begin_of_sentence|><|User|>\PH{请求}<|Assistant|></think>\PH{正文}<|end_of_sentence|><|User|>\PH{问题}<|Assistant|></think>
User       <|begin_of_sentence|><|User|>\PH{请求}\NL{}\NL{}\PH{正文}\NL{}\NL{}\PH{问题}<|Assistant|></think>
System     <|begin_of_sentence|><|User|>\PH{请求}<|System|>\PH{正文}<|User|>\PH{问题}<|Assistant|></think>
\end{Rendering}
\noindent 思考：
\begin{Rendering}
Assistant  <|begin_of_sentence|><|System|>Reasoning Effort: 75 (range 1-100, the higher the value, the more thorough the reasoning)\NL{}\NL{}<|User|>\PH{请求}<|Assistant|></think>\PH{正文}<|end_of_sentence|><|User|>\PH{问题}<|Assistant|><think>
User       <|begin_of_sentence|><|System|>Reasoning Effort: 75 (range 1-100, the higher the value, the more thorough the reasoning)\NL{}\NL{}<|User|>\PH{请求}\NL{}\NL{}\PH{正文}\NL{}\NL{}\PH{问题}<|Assistant|><think>
System     <|begin_of_sentence|><|System|>Reasoning Effort: 75 (range 1-100, the higher the value, the more thorough the reasoning)\NL{}\NL{}<|User|>\PH{请求}<|System|>\PH{正文}<|User|>\PH{问题}<|Assistant|><think>
\end{Rendering}
\paragraph{Ministral-3-14B-Instruct} （英文题目）
\noindent 不思考：
\begin{Rendering}
Start      <s>[SYSTEM_PROMPT]You are Ministral-3-14B-Instruct-2512, a Large Language Model (LLM) created by Mistral AI, a French startup headquartered in Paris.\NL{}You power an AI assistant called Le Chat.\NL{}Your knowledge base was last updated on 2023-10-01.\NL{}The current date is \LB{}today\RB{}.\NL{}\NL{}When you're not sure about some information or when the user's request requires up-to-date or specific data, you must use the available tools to fetch the information. Do not hesitate to use tools whenever they can provide a more accurate or complete response. If no relevant tools are available, then clearly state that you don't have the information and avoid making up anything.\NL{}If the user's question is not clear, ambiguous, or does not provide enough context for you to accurately answer the question, you do not try to answer it right away and you rather ask the user to clarify their request (e.g. "What are some good restaurants around me?" => "Where are you?" or "When is the next flight to Tokyo" => "Where do you travel from?").\NL{}You are always very attentive to dates, in particular you try to resolve dates (e.g. "yesterday" is \LB{}yesterday\RB{}) and when asked about information at specific dates, you discard information that is at another date.\NL{}You follow these instructions in all languages, and always respond to the user in the language they use or request.\NL{}Next sections describe the capabilities that you have.\NL{}\NL{}# WEB BROWSING INSTRUCTIONS\NL{}\NL{}You cannot perform any web search or access internet to open URLs, links etc. If it seems like the user is expecting you to do so, you clarify the situation and ask the user to copy paste the text directly in the chat.\NL{}\NL{}# MULTI-MODAL INSTRUCTIONS\NL{}\NL{}You have the ability to read images, but you cannot generate images. You also cannot read nor transcribe audio files or videos.\NL{}\NL{}# TOOL CALLING INSTRUCTIONS\NL{}\NL{}You may have access to tools that you can use to fetch information or perform actions. You must use these tools in the following situations:\NL{}\NL{}1. When the request requires up-to-date information.\NL{}2. When the request requires specific data that you do not have in your knowledge base.\NL{}3. When the request involves actions that you cannot perform without tools.\NL{}\NL{}Always prioritize using tools to provide the most accurate and helpful response. If tools are not available, inform the user that you cannot perform the requested action at the moment.[/SYSTEM_PROMPT]
Assistant  \PH{开头}[INST]\PH{请求}[/INST]\PH{正文}</s>[INST]\PH{问题}[/INST]
User       \PH{开头}[INST]\PH{请求}\NL{}\NL{}\PH{正文}\NL{}\NL{}\PH{问题}[/INST]
System     \PH{开头}[INST]\PH{请求}[/INST][SYSTEM_PROMPT]\PH{正文}[/SYSTEM_PROMPT][INST]\PH{问题}[/INST]
\end{Rendering}
\paragraph{Ministral-3-14B-Reasoning} （英文题目）
\noindent 思考：
\begin{Rendering}
Start      <s>[SYSTEM_PROMPT]# HOW YOU SHOULD THINK AND ANSWER\NL{}\NL{}First draft your thinking process (inner monologue) until you arrive at a response. Format your response using Markdown, and use LaTeX for any mathematical equations. Write both your thoughts and the response in the same language as the input.\NL{}\NL{}Your thinking process must follow the template below:[THINK]Your thoughts or/and draft, like working through an exercise on scratch paper. Be as casual and as long as you want until you are confident to generate the response to the user.[/THINK]Here, provide a self-contained response.[/SYSTEM_PROMPT]
Assistant  \PH{开头}[INST]\PH{请求}[/INST]\PH{正文}</s>[INST]\PH{问题}[/INST]
User       \PH{开头}[INST]\PH{请求}\NL{}\NL{}\PH{正文}\NL{}\NL{}\PH{问题}[/INST]
System     \PH{开头}[INST]\PH{请求}[/INST][SYSTEM_PROMPT]\PH{正文}[/SYSTEM_PROMPT][INST]\PH{问题}[/INST]
\end{Rendering}

在主探针下，DeepSeek-V4.1-Flash 的编码器不给开头的 Assistant 消息任何角色 token，Ministral 模板直接拒绝这样的对话，三者都把开头的 User 消息与问题合并。Ministral 各块中的\PH{开头}即上面自然顺序块中 Start 一行的文本。

\paragraph{DeepSeek-V4.1-Flash} （中文题目）
\noindent 不思考（未运行，仅作对照）：
\begin{Rendering}
Assistant  <|begin_of_sentence|>\PH{正文}<|end_of_sentence|><|User|>\PH{问题}<|Assistant|></think>
User       <|begin_of_sentence|><|User|>\PH{正文}\NL{}\NL{}\PH{问题}<|Assistant|></think>
System     <|begin_of_sentence|><|System|>\PH{正文}<|User|>\PH{问题}<|Assistant|></think>
\end{Rendering}
\noindent 思考（未运行，仅作对照）：
\begin{Rendering}
Assistant  <|begin_of_sentence|><|System|>Reasoning Effort: 75 (range 1-100, the higher the value, the more thorough the reasoning)\NL{}\NL{}\PH{正文}<|end_of_sentence|><|User|>\PH{问题}<|Assistant|><think>
User       <|begin_of_sentence|><|System|>Reasoning Effort: 75 (range 1-100, the higher the value, the more thorough the reasoning)\NL{}\NL{}<|User|>\PH{正文}\NL{}\NL{}\PH{问题}<|Assistant|><think>
System     <|begin_of_sentence|><|System|>Reasoning Effort: 75 (range 1-100, the higher the value, the more thorough the reasoning)\NL{}\NL{}\PH{正文}<|User|>\PH{问题}<|Assistant|><think>
\end{Rendering}
\paragraph{Ministral-3-14B-Instruct} （英文题目）
\noindent 不思考（未运行，仅作对照）：
\begin{Rendering}
Assistant  \Note{模板报错：Conversation must start with a user or system message, got assistant.}
User       \PH{开头}[INST]\PH{正文}\NL{}\NL{}\PH{问题}[/INST]
System     <s>[SYSTEM_PROMPT]\PH{正文}[/SYSTEM_PROMPT][INST]\PH{问题}[/INST]
\end{Rendering}
\paragraph{Ministral-3-14B-Reasoning} （英文题目）
\noindent 思考（未运行，仅作对照）：
\begin{Rendering}
Assistant  \Note{模板报错：Conversation must start with a user or system message, got assistant.}
User       \PH{开头}[INST]\PH{正文}\NL{}\NL{}\PH{问题}[/INST]
System     <s>[SYSTEM_PROMPT]\PH{正文}[/SYSTEM_PROMPT][INST]\PH{问题}[/INST]
\end{Rendering}
%>>app_templates
